\chapter{Ricci Flow as a Gradient Flow}

\section{Ricci Flow as a Gradient Flow}

\section{6.1 \quad Gradient of total scalar curvature and related functionals}

Given a compact manifold with a Riemannian metric $g$, we define the total scalar curvature by
\[
E(g)=\int_{\mathcal M} R\,dV.
\]

Let us consider the first variation of $E$ under an arbitrary change of metric.
As before, we write $h=\dfrac{\partial g}{\partial t}$. Using Propositions 2.3.9 and 2.3.12, and
keeping in mind that the integral of the divergence of a 1-form must be
zero, we compute
\[
\begin{aligned}
\frac{d}{dt}\int R\,dV
&=\int \frac{\partial R}{\partial t}\,dV+\int R\,\frac12(\tr h)\,dV\\
&=\int \bigl(-\langle \Ric,h\rangle+\delta^2h-\Delta(\tr h)\bigr)\,dV+\int \frac{R}{2}\langle g,h\rangle\,dV\\
&=\int \left\langle \frac{R}{2}g-\Ric,h\right\rangle\,dV.
\end{aligned}
\]

\begin{proposition}
\label{topping-gradient-flow-total-scalar-curvature-gradient}
\dcref{6.1.1}
\group{topping-gradient-flow}
\source{topping:page-0068,topping:page-0069}
The gradient of the total scalar curvature functional is
\[
\nabla E(g)=\frac{R}{2}g-\Ric.
\]
Moreover, if one formally evolves $g$ by the upward gradient flow $\partial_t g=\frac{R}{2}g-\Ric$, then the scalar curvature satisfies
\[
\frac{\partial R}{\partial t}
=-\left(\frac{n}{2}-1\right)\Delta R+|\Ric|^2-\frac12 R^2.
\]
\end{proposition}

\begin{proof}
The first variation computation above identifies the $L^2$-gradient of $E$ with $\frac{R}{2}g-\Ric$. Applying Proposition 2.3.9 together with the contracted Bianchi identity $\delta\Ric+\frac12 dR=0$ to the vector field $\partial_t g=\frac{R}{2}g-\Ric$ yields the stated evolution equation for $R$.
\end{proof}

One might first consider the flow of metrics up the gradient of $E$ -- that is,
the flow $\dfrac{\partial g}{\partial t}=\dfrac{R}{2}g-\Ric$.
Unfortunately, this flow cannot be expected to have any solutions starting at
an arbitrary metric $g_0$, in dimensions $n\geq 3$.
This is indicated by looking at how the scalar curvature would evolve under the
flow; by Proposition~2.3.9 and the identity $\delta\Ric+\frac{1}{2}dR=0$ we
would have
\[
\frac{\partial R}{\partial t}
=-\left(\frac{n}{2}-1\right)\Delta R+|\Ric|^2-\frac{1}{2}R^2,
\tag{6.1.1}
\]
which is a \emph{backwards heat equation}. Such equations do not typically admit
solutions, although this issue is complicated in this situation by the fact that
the Laplacian depends on the metric, and hence on the function $R$ for which
this is a PDE.

We should add that if only \emph{conformal} deformations of metric are allowed,
then the flow down the gradient is the so-called `Yamabe flow', also introduced
by Hamilton. (See for example [43].)

\section{6.2 \quad The $\mathcal{F}$-functional}

One of the early joys of the recent work of Perelman [31] is that Ricci flow
can in fact be formulated as a gradient flow. To see one such formulation, we
consider, for a fixed closed manifold $\mathcal M$, the following `Fisher
information' functional $\mathcal F$ on pairs $(g,f)$ where $g$ is a Riemannian
metric, and $f$ is a function $f:\mathcal M\to\mathbb{R}$:
\[
\mathcal{F}(g,f):=\int \left(R+|\nabla f|^2\right)e^{-f}\,dV.
\]

We need first to calculate the gradient of $\mathcal F$ under a smooth
variation of $g$ and $f$ with $\dfrac{\partial g}{\partial t}=h$ (as usual) and
$\dfrac{\partial f}{\partial t}=k$, for some function $k:\mathcal M\to\mathbb{R}$.

\begin{proposition}
\label{topping-gradient-flow-f-functional-first-variation}
\dcref{6.2.1}
\group{topping-gradient-flow}
\source{topping:page-0069,topping:page-0070}
Let $\mathcal F(g,f)=\int (R+|\nabla f|^2)e^{-f}\,dV$. For a variation
\[
\frac{\partial g}{\partial t}=h, \qquad \frac{\partial f}{\partial t}=k,
\]
its first variation is
\[
\frac{d}{dt}\mathcal{F}(g,f)
=\int \langle-\Ric-\Hess(f),h\rangle e^{-f}\,dV
+\int \left(2\Delta f-|\nabla f|^2+R\right)\left(\frac{1}{2}\tr h-k\right)e^{-f}\,dV.
\]
\end{proposition}

\begin{proof}
Differentiating $|\nabla f|^2$ gives
\[
\frac{\partial}{\partial t}|\nabla f|^2
=-h(\nabla f,\nabla f)+2\langle\nabla k,\nabla f\rangle.
\]
Using Propositions 2.3.9 and 2.3.12, we obtain
\[
\frac{d}{dt}\mathcal F(g,f)
=\int \bigl[-h(\nabla f,\nabla f)+2\langle\nabla k,\nabla f\rangle-\langle\Ric,h\rangle+\delta^2h-\Delta(\tr h)\bigr]e^{-f}\,dV
\]
\[
+\int (R+|\nabla f|^2)\bigl[-k+\tfrac12\tr h\bigr]e^{-f}\,dV.
\]
Integrating by parts the $2\langle\nabla k,\nabla f\rangle$, $\delta^2h$, and
$\Delta(\tr h)$ terms and then collecting coefficients yields the stated formula.
\end{proof}

\section{6.3 \quad The heat operator and its conjugate}

We pause to define some notation. We denote the heat operator, acting on
functions $f : \mathcal{M} \times [0,T] \to \mathbb{R}$ by
\[
  \Box := \frac{\partial}{\partial t} - \Delta.
\]
We also need to define
\[
  \Box^{*} := - \frac{\partial}{\partial t} - \Delta + R,
\]
which is conjugate to $\Box$ in the following sense.

\begin{proposition}
\label{topping-gradient-flow-conjugate-heat-operator}
\dcref{6.3.1}
\group{topping-gradient-flow}
\source{topping:page-0071}
If $g(t)$ is a Ricci flow for $t \in [0,T]$, and $v,w : \mathcal{M} \times
[0,T] \to \mathbb{R}$ are smooth, then
\[
  \int_0^T \left( \int_{\mathcal{M}} (\Box v) w\, dV \right) dt
  = \left[\int_{\mathcal{M}} v w\, dV\right]_0^T
  + \int_0^T \left( \int_{\mathcal{M}} v(\Box^{*} w)\, dV \right) dt.
\]
\end{proposition}

\begin{proof}
Differentiate $\int vw\,dV$ in time and use the Ricci-flow volume variation
formula together with integration by parts in space. Rewriting the resulting
terms using $\Box$ and $\Box^*$ gives the claimed identity.
\end{proof}

\begin{remark}
\label{topping-gradient-flow-conjugate-heat-special-case}
\dcref{6.3.2}
\group{topping-gradient-flow}
\source{topping:page-0071}
The special case $v\equiv 1$ gives
\[
  \frac{d}{dt} \int w\, dV = - \int \Box^{*} w\, dV.
\]
\end{remark}

\section{6.4 \quad A gradient flow formulation}

Let us now return to the discussion of Section 6.2, specialising to
variations of $g$ and $f$ which preserve the distorted volume form
$e^{-f} dV$. From Proposition 2.3.12, we see that
\[
  0 = \frac{\partial}{\partial t}(e^{-f} dV)
  = \left(\frac{1}{2}\tr \frac{\partial g}{\partial t}
  - \frac{\partial f}{\partial t}\right)e^{-f} dV,
\]
and so the evolution of $f$ is determined by
\[
  \frac{\partial f}{\partial t}
  = \frac{1}{2}\tr \frac{\partial g}{\partial t}.
\]
The evolution of $\mathcal F$ then proceeds according to
\[
  \frac{d}{dt}\mathcal{F}(g,f)
  = \int \langle -\Ric - \Hess(f), \frac{\partial g}{\partial t}\rangle e^{-f} dV.
\]
From here, we can see that if we had a solution to the flow
\[
  \frac{\partial g}{\partial t} = -2(\Ric + \Hess(f)),
  \tag{6.4.1}
\]
and the corresponding
\[
  \frac{\partial f}{\partial t} = -R - \Delta f,
  \tag{6.4.2}
\]
then $\mathcal{F}$ would evolve under
\[
  \frac{d}{dt}\mathcal{F}(g,f)
  = 2 \int \lvert \Ric + \Hess(f)\rvert^2 e^{-f} dV \geq 0.
  \tag{6.4.3}
\]
In this case, defining
\[
  \omega := e^{-f} dV,
  \tag{6.4.4}
\]
which would now be constant in time, we could view $g$ as a gradient flow for
the functional $g \mapsto \mathcal{F}(g,f)$, where $f$ is determined from
$\omega$ and $g$ by (6.4.4).

The coupled system (6.4.1) and (6.4.2) looks at first glance to be somewhat
alarming from a PDE point of view. Whilst (6.4.1) resembles the Ricci flow,
which can be solved forwards in time for given initial data, (6.4.2) is a
backwards heat equation and would not admit a solution for most prescribed
initial data.

The trick to obtaining a solution of this system is to show that it is somehow
equivalent to the decoupled system of equations
\[
  \frac{\partial g}{\partial t} = -2\Ric,
  \tag{6.4.5}
\]
\[
  \frac{\partial f}{\partial t} = -\Delta f + \lvert \nabla f\rvert^2 - R,
  \tag{6.4.6}
\]

By the short-time existence of Theorem~5.2.1, (6.4.5) admits a smooth solution, for given initial data $g_0$, on some time interval $[0,T]$. As before, (6.4.6) is a backwards heat equation, and we cannot expect there to exist a solution achieving an arbitrary initial $f(0)$. Instead, for the given smooth $g(t)$, we prescribe \emph{final} data $f(T)$ and solve backwards in time. Such a solution exists all the way down to time $t=0$, for each smooth $f(T)$, as we can see by changing variables to

\[
u(t) := e^{-f(t)},
\tag{6.4.7}
\]

which satisfies the simple linear equation

\[
\Box^{*}u = 0,
\tag{6.4.8}
\]

which admits a positive solution backwards to time $t=0$.

Given these $g(t)$ and $f(t)$, we set $X(t) := -\nabla f$ (which generates a unique family of diffeomorphisms $\psi_t$ with $\psi_0 = \id$, say) so that $\mathcal L_X g = -2\Hess(f)$ by (2.3.9). We fix $\sigma(t) \equiv 1$. Now when we define $\hat g(t)$ by (1.2.2), it will evolve, by (1.2.3), according to

\[
\frac{\partial \hat g}{\partial t}
= \psi_t^{*}\bigl[-2\Ric(g) - 2\Hess_g(f)\bigr].
\]

where $\Hess_g(f)$ is the Hessian of $f$ with respect to the metric $g$. Keeping in mind that $\psi_t : (\mathcal M,\hat g) \to (\mathcal M,g)$ is an isometry, we may then write $\hat f := f \circ \psi_t$ to give

\[
\frac{\partial \hat g}{\partial t}
= -2\bigl(\Ric(\hat g) + \Hess_{\hat g}(\hat f)\bigr).
\tag{6.4.9}
\]

The evolution of $\hat f$ is then found by the chain rule:

\[
\begin{aligned}
\frac{\partial \hat f}{\partial t}(x,t)
&= \frac{\partial f}{\partial t}(\psi_t(x),t) + X(f)(\psi_t(x),t) \\
&= \bigl[(-\Delta_g f + \lvert \nabla f\rvert^2 - R_g) - \lvert \nabla f\rvert^2\bigr](\psi_t(x),t) \\
&= \bigl[-\Delta_{\hat g}\hat f - R_{\hat g}\bigr](x,t),
\end{aligned}
\]

where it has been necessary again to use subscripts to denote the metric being used. Thus we have a solution to

\[
\frac{\partial \hat f}{\partial t} = -\Delta_{\hat g}\hat f - R_{\hat g}
\tag{6.4.10}
\]

coupled with (6.4.9), as desired. We have shown that solutions of (6.4.1) and (6.4.2) may be generated by pulling back solutions of (6.4.5) and (6.4.6)
by an appropriate time-dependent diffeomorphism. One may also work the other way, finding solutions of (6.4.5) and (6.4.6) given solutions of (6.4.1) and (6.4.2).

Note that pulling back $g$ and $f$ by the same diffeomorphism leaves $\mathcal F$ invariant. In particular, $\mathcal F$ is monotonically increasing as $g$ and $f$ evolve under (6.4.5) and (6.4.6).

Let us take stock of what we have shown, in the following theorem.

\begin{theorem}
\label{topping-gradient-flow-gradient-flow-formulation}
\dcref{6.4.1}
\group{topping-gradient-flow}
\source{topping:page-0074}
\uses{topping-gradient-flow-f-functional-first-variation,topping-gradient-flow-conjugate-heat-operator,topping:variation-volume-form}
For an arbitrary $g_0$, on a closed manifold $\M$, there exist $T>0$, and an $n$-form $\omega>0$ which is neither arbitrary nor unique, such that there exists a solution on the time interval $[0,T]$ to the (upwards) gradient flow for the functional
\[
\mathcal F^\omega(g)=\int (R+\lvert \nabla f\rvert^2)\,\omega,
\]
where $f$ is determined by the constraint $\omega=e^{-f}dV$, and we compute the gradient with respect to the inner product
\[
\langle g,h\rangle_\omega=\frac12 \int \langle g,h\rangle\,\omega.
\]
\end{theorem}

\begin{proof}
The coupled system for $(g,f)$ is related to the decoupled system by pulling back along the diffeomorphisms generated by $X=-\nabla f$. The distorted volume form $\omega=e^{-f}dV$ is fixed along the coupled system, and the first variation formula for $\mathcal F$ shows that the resulting flow is the upward gradient flow for $\mathcal F^\omega$.
\end{proof}

We are free to take $T$ as large as we like, provided that the Ricci flow starting at $g_0$ exists for $t\in[0,T]$, although $\omega$ will depend on the $T$ chosen.

This gradient flow is just the Ricci flow, modified by a time-dependent diffeomorphism. That is, there exists a smooth family of diffeomorphisms $\phi_t:\M\to\M$ such that $\phi_t^*(g(t))$ is the Ricci flow starting at $g_0$.

\begin{remark}
\label{topping-gradient-flow-volume-preserving-remark}
\dcref{6.4.2}
\group{topping-gradient-flow}
\source{topping:page-0074}
In the construction above, the measure $\omega$ can be described either from the initial data $(\hat g(0),\hat f^{(0)})$ or from the terminal data $(\hat g(T),f(T))$ after pulling back by $\psi_T$.
\end{remark}

\begin{remark}
\label{topping-gradient-flow-steady-soliton-criterion}
\dcref{6.4.3}
\group{topping-gradient-flow}
\source{topping:page-0074}
If the time derivative of $\mathcal F$ vanishes along the system (6.4.1)--(6.4.2), then $\Ric+\Hess(f)=0$. On a compact manifold this forces the associated Ricci flow to be Ricci flat, so the flow is a steady gradient soliton.
\end{remark}

\section{6.5 \quad The classical entropy}

So far in Chapter 6, our exposition has been an expansion of the original
comments of Perelman [31]. It is enlightening also to emphasise the function
$u$ from (6.4.7) rather than $f$. If we specify $u(T)$ as a probability density
at time $T$, and let it evolve backwards in time under (6.4.8), then $u(t)$
represents the probability density of a particle evolving under Brownian
motion, backwards in time, on the evolving manifold. A basic building block
for constructing interesting functionals for Ricci flow is then the classical,
or `Boltzmann-Shannon' entropy
\[
  N = - \int_{\mathcal{M}} u \ln u\, dV,
\]
as considered in thermodynamics, information theory (Boltzmann, Shan-
non, Turing etc.) and elsewhere, which gives a measure of the concentration
of the probability density. (Note that Brownian diffusion on a fixed mani-
fold represents the gradient flow for this entropy on the space of probability
densities endowed with the Wasserstein metric -- see the discussion and ref-
erences in [42].) In this framework, the information functional $\mathcal{F}$
coincides with the time derivative of $N$, up to a sign (see also [30]). By (2.5.7),
\[
  -\frac{dN}{dt}
  = \int \left[ (1 + \ln u)\frac{\partial u}{\partial t} - Ru\ln u \right] dV
  = - \int \left[ (1 + \ln u)(\Delta u - Ru) + Ru\ln u \right] dV
  \tag{6.5.1}
\]
\[
  = \int \left( \frac{|\nabla u|^2}{u} + Ru \right) dV
\]
\[
  = \mathcal{F}.
\]

It is sometimes convenient to take a renormalised version of the classical
entropy
\[
  \widetilde{N} = N - \frac{n}{2}\left( 1 + \ln 4\pi (T - t) \right).
  \tag{6.5.2}
\]

With this adjustment, one can check that $\widetilde{N} \to 0$ as $t \uparrow T$
when one takes $u$ to be a solution of (6.4.8) which is a delta function at
$t = T$. (In other words, $u$ is a heat kernel for the conjugate heat equation --
the probability density of the backwards Brownian path of a particle whose
position is known at time $t = T$.) Indeed, this functional applied to the
classical fundamental solution of the heat equation on Euclidean space would
be zero. One could check that by virtue of the Harnack inequality of Li-Yau [28],
the functional $\widetilde{N}$ applied to a (backwards) heat kernel on a fixed
manifold of positive Ricci curvature must be decreasing in $T - t$. For Ricci flow,
this is true without any curvature hypothesis:

\begin{lemma}
\label{topping-gradient-flow-classical-entropy-monotone}
\dcref{6.5.1}
\group{topping-gradient-flow}
\source{topping:page-0076}
\uses{topping-gradient-flow-gradient-flow-formulation}
If $g(t)$ is a Ricci flow on a closed manifold $M$, for $t\in[0,T]$, and $u:M\times[0,T]\to(0,\infty)$ is a solution to $\Box^{*}u=0$, then the functional $\widetilde N$ defined by (6.5.2) is monotonically increasing in $t$.
\end{lemma}

\begin{proof}
As observed above, we have
\[
-\frac{d\widetilde N}{dt}=\mathcal F-\frac{n}{2(T-t)},
\]
so we are reduced to proving that
\begin{equation}\tag{6.5.3}
\mathcal F\le \frac{n}{2(T-t)}.
\end{equation}
The functional $\mathcal F$ evaluated as here on the Ricci flow $g(t)$ and the density $u$ (or its corresponding $f:=-\ln u$) is equal to $\mathcal F$ evaluated on the $\hat g$ and $\hat f$ of the previous section. We may then sharpen the monotonicity of $\mathcal F$, as expressed in (6.4.3) to
\begin{align}
\frac{d\mathcal F}{dt}
&=2\int \lvert \Ric+\Hess(f)\rvert^2 e^{-f}\,dV
\ge \frac{2}{n}\int \lvert R+\Delta f\rvert^2 e^{-f}\,dV \tag{6.5.4}\\
&\ge \frac{2}{n}\left(\int (R+\Delta f)e^{-f}\,dV\right)^2
=\frac{2}{n}\mathcal F^2,
\end{align}
where we have estimated the norm of the symmetric tensor $\Ric+\Hess(f)$ in terms of its trace, appealed to Cauchy-Schwarz, and integrated by parts, as in [31]. If ever (6.5.3) were violated, then the differential inequality (6.5.4) would force $\mathcal F(t)$ to blow up prior to the time $T$ until which we are assuming its existence.
\hfill$\square$
\end{proof}

Essentially, this lemma says that the evolution of a manifold under Ricci flow will always balance with backwards Brownian motion to prevent diffusion faster than that on Euclidean space.

Currently the most useful entropy functional for Ricci flow is Perelman's $\mathcal W$-entropy which we shall meet in Chapter 8, and this too will arise from considering the entropy $N$.

\section{6.6 \quad The zeroth eigenvalue of \texorpdfstring{$-4\Delta + R$}{-4 Delta + R}}

It is also useful to rewrite the functional $\mathcal{F}$ to act on pairs $(g,\phi)$ where
$\phi := e^{-f/2}$. In this case, one can easily check that
\[
\mathcal{F} = \int \bigl(4|\nabla\phi|^2 + R\phi^2\bigr)\,dV,
\]
and hence that the zeroth eigenvalue of the operator $-4\Delta + R$ is given by
\[
\lambda = \lambda(g) = \inf\left\{\mathcal{F}(g,\phi)\ \middle|\ \int \phi^2\,dV = 1\right\},
\]
where we are using the same notation $\mathcal{F}$ even after the change of variables.
By applying the direct method and elliptic regularity theory (see [15, §8.12])
we know that this infimum is attained by some smooth $\phi_{\min}$ (for given
smooth $g$ on a closed manifold) and also that after negating if necessary, we
have $\phi_{\min} > 0$. (Note that $|\phi_{\min}|$ will also be a minimiser, and that by the
Harnack inequality -- see for example [15, Theorem 8.20] -- we must then
have $|\phi_{\min}| > 0$.)

One consequence of this positivity is that the minimiser $\phi_{\min}$ must be unique
(up to a sign) or we would be able to take an appropriate linear combination
of two distinct minimisers to give a new minimiser which would no longer
be nonzero.

Another consequence is that we can write $\phi_{\min} = e^{-f_{\min}/2}$ for some function
$f_{\min} : \mathcal{M} \to \mathbb{R}$, and hence we see that
\begin{equation}
\lambda = \min\left\{\mathcal{F}(g,f)\ \middle|\ \int e^{-f}\,dV = 1\right\}.
\tag{6.6.1}
\end{equation}

\begin{theorem}
\label{topping-gradient-flow-z0-eigenvalue-monotone}
\dcref{6.6.1}
\group{topping-gradient-flow}
\source{topping:page-0077,topping:page-0078}
\uses{topping-gradient-flow-gradient-flow-formulation}
If $g(t)$ is a Ricci flow on a closed manifold $\mathcal{M}$, for $t \in [0,T]$, then the eigenvalue $\lambda(g(t))$ is a weakly increasing function on $[0,T]$.
\end{theorem}

\begin{proof}
Pick arbitrary times $a$ and $b$ with $0 \le a < b \le T$, and define a
function $f_b$ to be the minimiser in (6.6.1) corresponding to the metric $g(b)$.
Now if we define a time-dependent function $f(t)$ for $t \in [a,b]$ to solve (6.4.6)
with final condition $f(b) = f_b$ (there is a unique smooth solution as in
Section 6.4) then by the monotonicity of Section 6.4, we see that
\[
\lambda(g(a)) \le \mathcal{F}(g(a),f(a)) \le \mathcal{F}(g(b),f(b)) = \lambda(g(b)).
\]
\end{proof}

The monotonicity of such an eigenvalue may be considered -- when $\lambda > 0$ --
to be equivalent to the improvement of the best constant in an appropriate
Poincar\'e inequality. In Chapter 8, we will see that the best constant in a
certain log-Sobolev inequality is also improving.
